% Variable fields such as your name, title of the thesis atc. are set in this file.
% The file is SHARED between the thesis text and the presentation --- no need to set anything twice.

\usepackage[
  english,			% English study program
%%% Choose one of the thesis types
  %semestral,		%	semestral project (default)
  %bachelor,			%	Bachelor's thesis
  master,			% Master's thesis
  %treatise,			% Treatise on doctoral thesis
  %doctoral,				% Doctoral thesis
%%% Choose one of the options:
%  left,				% Equations and captions will be aligned to the left
  center,			% Equations and captions will be centered (default)
%
  %electronic       % Electronically published version of the thesis; page margins do not alternate with odd- and even-numbered pages
]{thesis}   % Package for typesetting theses


%%% First and last name of the thesis author, following the scheme
% [titles in front]{FirstName}{LastName}[titles after]
% If the author does not have a title in front/after, just delete the entire string '[...]'
\author[Bc.]{Vladislav}{Válek}

%%% Brno University of Technology identification number of author (BUT ID)
\butid{4738967}

%%% First and last name of the advisor
% [titles in front]{FirstName}{LastName}[titles after]
% If the author does not have a title in front/after, just delete the entire string '[...]'
% [titles in front]{FirstName}{LastName}[titles after]
\advisor[Prof. Dr.]{Dirk}{Koch}

%%% First and last name of the opponent
% [titles in front]{FirstName}{LastName}[titles after]
% Makes use only in the presentation for defense;
% in case you do not want the opponent to be shown on the title slide, just comment out the command
% The opponent is not shown in case of the semestral thesis (no opponent exists actually at that time)
% In the case of a doctoral thesis, two to three oponents are typically involved. In such a case, if you want to have them on the title slide, please go to the definition of "VUT title page" in the thesis.sty file, uncomment and adapt their names.
\opponent[Prof. Dr.]{Holger}{Fröning}

%%% Thesis title
%  In the case of a very long thesis title, it may happen that it does not fit
%  into the slot in the footbar of the slides. You may use the command
%  \def\insertshorttitle{Shortened th.\ title}
%  where the shortened version of the title appears as the parameter.
%  If you do not want to shorten the title, you will have to redefine how the slide footbar
%  is generated, see: https://bit.ly/3EJTp5A
\title{Shared CPU-FPGA Filesystem Access on fast NVMe Devices}

%%% Study program/specialization
\specialization{Computer Engineering} %Teleinformatika

%%% Department
%\department{Department of Control and Instrumentation}  %Ústav automatizace a měřicí techniky
%\department{Department of Biomedical Engineering}       %Ústav biomedicínského inženýrství
%\department{Department of Electrical Power Engineering} %Ústav elektroenergetiky
%\department{Department of Electrical and Electronic Technology}   %Ústav elektrotechnologie
%\department{Department of Physics}                      %Ústav fyziky
%\department{Department of Foreign Languages}            %Ústav jazyků
%\department{Department of Mathematics}                  %Ústav matematiky
%\department{Department of Microelectronics}             %Ústav mikroelektroniky
%\department{Department of Radio Electronics}            %Ústav radioelektroniky
%\department{Department of Theoretical and Experimental Electrical Engineering}  %Ústav teoretické a experimentální elektrotechniky
% \department{Department of Telecommunications}           %Ústav telekomunikací
%\department{Department of Power Electrical and Electronic Engineering}   %Ústav výkonové elektrotechniky a elektroniky
\department{Institute of Computer Engineering (ZITI)}

%%% Faculty
%\faculty{Faculty of Architecture}   %Fakulta architektury
% \faculty{Faculty of Electrical Engineering and~Communication}   %Fakulta elektrotechniky a~komunikačních technologií
%\faculty{Faculty of Chemistry}   %Fakulta chemická
%\faculty{Faculty of Information Technology}   %Fakulta informačních technologií
%\faculty{Faculty of Business and Management}   $Fakulta podnikatelská
%\faculty{Faculty of Civil Engineering}   %Fakulta stavební
%\faculty{Faculty of Mechanical Engineering}   %Fakulta strojního inženýrství
%\faculty{Faculty of Fine Arts}   %Fakulta výtvarných umění
\faculty{Faculty of Engineering Sciences}
%
%Logotype selection (in square brackets short logo, in curly brackets full logo):
\facultylogo[logo/FEEC_abbreviation_color_PANTONE_EN]{logo/UTKO_color_PANTONE_EN}


%%% Graduate year (typically the calendar year of the defense)
\graduateyear{2026}
%%% Academic year (typically the year of solution of the thesis in the format n/n+1)
\academicyear{2025/26}
% Date of the defense (makes use only in the presentation slides)
\date{30-March-2026}

%%% Place of the defense
\city{Heidelberg}

%%% Abstract
\abstract{%
This thesis presents an architecture for shared CPU-FPGA filesystem access to NVMe storage using
PCI Express peer-to-peer communication. The goal is to reduce host-memory staging and software
overhead by enabling direct accelerator interaction with an NVMe device while preserving standard
host-side filesystem usage. The implementation combines hardware and software components: an FPGA IP
core (DMA Iuventus) integrated with PCIe endpoint interfaces processes NVMe-related traffic, manages
queue interactions, and issues read/write commands; a userspace software stack based on NDK-SW and
SPDK, together with Linux ublk, provides control-plane integration and exposes a mountable block
device. The complete flow is validated with an exFAT filesystem and demonstrates shared access from
both CPU and FPGA contexts. Experimental evaluation on an AMD Alveo U55C platform with an NVMe SSD
reports resource utilization and request latency for multiple transfer sizes and access patterns. The
results confirm functional operation of the proposed approach and show distinct latency
characteristics for read and write paths, with sublinear read-latency growth for larger transfers.
The work demonstrates practical feasibility of host-bypass storage acceleration and provides a
foundation for future optimization of synchronization, throughput, and multi-device deployment
scenarios.
}

%%% Keywords
\keywrds{%
PCIe, DMA, host-bypassing, peer-to-peer, SPDK, NVMe, FPGA, userspace, ublk, bdev
}

%%% Thanks and acknowledgement
\acknowledgement{%
I would like to thank my supervisor, Prof.\ Dr.\ Dirk Koch for his valuable feedback when writing
this thesis as well as in my research work at the Heidelberg University. I would also like to thank
my colleague, Dr.\ Riadh Ben Abdelhamid, for providing necessary user feedback for the current
course of research. Further, I give many thanks to my colleagues at CESNET, a.l.e., Ing.\ Jakub
Cabal and Ing.\ Martin Špinler, for mentoring me in the development of PCIe-compliant IPs and
software. Finally, I thank my mom, Ing.\ Marta Válková, for supporting me throughout my studies.
}%
